% Generated draft: requires mathematical and source review.
\documentclass[11pt]{article}
\usepackage[margin=25mm]{geometry}
\usepackage{fontspec}
\setmainfont{Latin Modern Roman}
\usepackage{amsmath,amssymb,mathtools,unicode-math}
\setmathfont{Latin Modern Math}
\usepackage{enumitem,needspace,xurl,hyperref}
\setlength{\parindent}{0pt}
\setlength{\parskip}{5pt}
\setlength{\emergencystretch}{3em}
\newcommand{\N}{\mathbb{N}}
\newcommand{\Z}{\mathbb{Z}}
\newcommand{\Q}{\mathbb{Q}}
\newcommand{\R}{\mathbb{R}}
\newcommand{\C}{\mathbb{C}}
\newcommand{\F}{\mathbb{F}}
\newcommand{\A}{\mathbb{A}}
\newcommand{\PP}{\mathbb{P}}
\newcommand{\E}{\mathbb{E}}
\newcommand{\eps}{\varepsilon}
\newcommand{\dd}{\,\mathrm{d}}
\begin{document}
\section*{UnsolvedMath ID: 30005265}
\textbf{Source problem code:} catalogue entry
\subsection{Definitions and mathematical statement}
A Dirichlet character modulo a prime \(q\) is a group homomorphism \(\chi:(\Z/q\Z)^\times\to\C^\times\), extended periodically to integers by setting \(\chi(n)=0\) when \(q\mid n\). It is primitive if its conductor is \(q\). Its order \(\operatorname{ord}(\chi)\) is the least positive integer \(d\) such that \(\chi(a)^d=1\) for every \(a\in(\Z/q\Z)^\times\). For a nonprincipal character, define
\[
\ell(\chi)=\min\{n\in\Z_{\geq1}:\chi(n)\notin\{0,1\}\}.
\]
This minimum exists because the character is nonprincipal.

The supplied problem asks for general conditions on primitive characters modulo primes, with orders tending to infinity, that ensure
\[
\ell(\chi)=O_{\varepsilon}(q^{\varepsilon})
\qquad\text{for every }\varepsilon>0,
\]
and also ensure that \(\chi\) takes ``varied values'' below \(q^{\varepsilon}\).

To make the limiting regime explicit, one may consider a family \((q_j,\chi_j)_{j\geq1}\), where each \(q_j\) is prime, each \(\chi_j\) is primitive modulo \(q_j\), and \(\operatorname{ord}(\chi_j)\to\infty\) as \(j\to\infty\). The requested small-value bound then has the schematic quantifier form
\[
\forall\varepsilon>0\ \exists C_{\varepsilon}>0\ \exists j_0\
\forall j\geq j_0:\quad
\ell(\chi_j)\leq C_{\varepsilon}q_j^{\varepsilon}.
\]
The problem is to identify hypotheses under which such a bound holds. The supplied record does not specify those hypotheses, the required growth of the character orders relative to the moduli, or whether the constant and threshold must be uniform across all admissible families. The display therefore describes the requested bound within a fixed family and does not resolve the intended uniformity.

For clarity, an auxiliary set recording distinct nonzero values is
\[
V_{\chi}(X)=\{\chi(n):n\in\Z_{\geq1},\ n<X,\ q\nmid n\}.
\]
The supplied phrase ``varied values'' does not specify a lower bound for \(\lvert V_{\chi}(q^{\varepsilon})\rvert\), an equidistribution condition, or another quantitative property. It also does not say whether zero values are relevant or explicitly settle the endpoint convention. Thus the second requested conclusion cannot be stated as a definite mathematical predicate from the supplied evidence. No particular interpretation of that conclusion is asserted here.
\subsection{Short English statement}
Find conditions on primitive Dirichlet characters of prime modulus, with orders tending to infinity, that force the least positive integer with character value outside \(\{0,1\}\) to be bounded by every positive power of the modulus, and force diversity of small character values. The required uniformity and the meaning of value diversity remain unspecified.
\subsection{Sources}
\begin{itemize}
\item Local UnsolvedMath registry, supplied record with permanent ID \(30005265\), title ``Vinogradov-Type Bounds for Large-Order Dirichlet Characters,'' exact source code \texttt{catalogue entry}; collection: Oberwolfach Reports Open Problems.
\item Supplied report citation: \emph{Analytic Number Theory} (2023), Oberwolfach Reports, pp.\ 2895--2960, DOI \texttt{10.4171/owr/2022/50}. \url{https://doi.org/10.4171/owr/2022/50}. The record separately supplies proposed year 2022; these year fields are retained without reconciliation.
\item Attribution supplied in the original-statement field: Mangerel, ``Large order Dirichlet characters and an analogue of a conjecture of Vinogradov.'' No proposer is identified separately. Supplied literature reference labelled ``Mangerel'': \url{https://doi.org/10.1112/jlms.12841}.
\item Supplied literature reference labelled ``Oberwolfach Report 50'': \url{https://ems.press/content/serial-article-files/46986}.
\item Literature assessment dated exactly 2026-08-22: the supplied curation labels the problem open and states that the cited work proves substantial large-order character estimates while the extracted analogue of Vinogradov's conjectural problem remains unsettled in general. This is supplied dated curation derived from the OWR material, not an independently verified literature assessment. The source requests independent verification of time-sensitive status and cited claims.
\item The registry labels the statement \texttt{corrected\_verified} and states that it was reconstructed and checked against the official Oberwolfach report. These are supplied provenance claims, not verification performed for this draft.
\end{itemize}
\end{document}
